\documentclass[12pt,a4paper]{report}
\usepackage[utf8]{inputenc}
\usepackage{setspace}
\usepackage{geometry}
\usepackage{graphicx}
\usepackage{pdfpages}
\usepackage{hyperref}
\usepackage{enumitem}
\usepackage{float}
\usepackage{cleveref}
\usepackage{booktabs}
\geometry{margin=1in}
\onehalfspacing
\usepackage{newunicodechar}
\newunicodechar{ }{\,} % U+202F NARROW NO-BREAK SPACE

\renewcommand{\bibname}{References}

\begin{document}

%=========================
% COVER PAGE
\begin{titlepage}
    \centering
    \vspace*{2cm}
    
    % University Logo Placeholder
    %\includegraphics[width=0.4\textwidth]{university_logo.png}\par
    %\vspace{1cm}
    
    {\large\bfseries COMP2043.GRP Final Group Report\par}
    \vspace{1.5cm}
    
    % University slogan placeholder
    {\Large \textbf{Integrating Artificial Intelligence (AI) and Virtual Production (VP) workflows in implementing and testing a framework for Chinese short films with international appeal}\par}
    \vspace{1cm}
    
    % Team Number
    {\Large \textbf{Team 17}\par}
    \vspace{0.4cm}
    
    % Team Crew
    \begin{spacing}{1.2}
    \large\textbf{Qingyang Liu (sayql5) 20617232} \\
    \large\textbf{Liangli Wang (scylw2) 20616381} \\
    \large\textbf{Youyang Wu (scyyw26) 20616341} \\
    \large\textbf{Houpu Song (scyhs6) 20617167} \\
    \large\textbf{Haoyi Jiang (scyhj6) 20617498} \\
    \end{spacing}
    \vspace{1cm}
    
    % Supervisors
    {\Large Supervised by\par}
    \begin{spacing}{1.5}
    \large\textbf{Prof. Dave TOWEY} \\
    \large\textbf{Dr. Nazia ANWAR} \\
    \end{spacing}
    
    \vfill
    {\large University of Nottingham Ningbo China \par}
    \vspace{0.15cm}
    {\large \today\par}
\end{titlepage}


%=========================
% Abbreviations
\newpage
\addcontentsline{toc}{section}{Abbreviations}
\begin{spacing}{2}
    \renewcommand{\arraystretch}{1.5}
    \begin{center}
    \begin{tabular*}{\textwidth}{@{\extracolsep{\fill}}ll@{}}
        \toprule
        \large{\textbf{Abbreviation}} & \large{\textbf{Full Term}} \\
        \midrule
        \textbf{ARRI} & Arnold \& Richter Cine Technik \\
        \textbf{DoP} & Director of Photography \\
        \textbf{ICVFX} & In-Camera Visual Effects \\
        \textbf{LOD} & Level of Detail \\
        \textbf{SRS} & Software Requirements Specification \\
        \textbf{SE} & Software Engineering \\
        \textbf{TET} & Technical Equipment Team \\
        \textbf{UE / UE5} & Unreal Engine (Version 5) \\
        \textbf{VAD} & Virtual Art Department \\
        \textbf{VP} & Virtual Production \\
        \bottomrule
    \end{tabular*}
    \end{center}
    \renewcommand{\arraystretch}{1}
\end{spacing}


%=========================
% TABLE OF CONTENTS
\newpage
\begin{spacing}{1.8}
    \tableofcontents
\end{spacing}
\newpage


%=========================
% CHAPTER 1: INTRODUCTION
\chapter{Introduction}

\section{Project Context}

\section{Core Objectives}

\section{Overview}


%=========================
% CHAPTER 2: BACKGROUND RESEARCH
\chapter{Background Research}

\section{VP Pipelines \& ICVFX}
Virtual Production (VP) induces a paradigm shift in traditional filmmaking pipelines by replacing physical sets with real-time digital environments, thereby shifting a substantial amount of post-production decision-making to the pre-production phase \cite{kadner2019vp}. In-Camera Visual Effects (ICVFX) serves as the technical hub of this pipeline, requiring ultra-low latency hardware-level synchronization among camera tracking systems, real-time rendering engines, and LED display matrices. Particularly in VP projects focused on small-scale interior scenes, the spatial distance between the virtual environment and physical props is heavily compressed. Under these constraints, the visual realism of the system depends not merely on the polygon count of digital assets, but strictly relies on the precise alignment of physical props, dynamic lighting, and spatial perspective during the on-set formal shoot \cite{swords2024emergence}.

\section{AI in VP}
The intervention of Artificial Intelligence significantly enhances VP production efficiency. During the asset construction phase, generative AI tools can rapidly convert natural language or 2D images into 3D meshes equipped with Physically Based Rendering (PBR) materials. Concurrently, neural rendering technologies (e.g., NVIDIA DLSS) substantially reduce the GPU computational load required for rendering high-fidelity scenes via deep learning super sampling \cite{alarcon2020dlss}. However, regarding the engineering bottleneck of unaligned virtual and physical spaces on-set, existing AI post-processing tools still suffer from severe temporal and spatial inconsistencies when attempting to repair complex physical perspectives and depth-of-field occlusions. This objectively highlights the necessity of developing in-engine, real-time alignment support tools.

\section{Case Studies and Industry Practices}
Current industry practices provide feasibility demonstrations for intelligent VP. In industrial-standard productions, the series \textit{The Mandalorian} validated the physical optical effectiveness of LED wall ``lighting from the source,'' which provides interactive illumination directly to the physical set from the digital environment, thereby eliminating green screen spill \cite{kadner2019vp}. Regarding agile engineering collaboration, the 5-person visual effects team for the film \textit{Everything Everywhere All At Once} successfully delivered over 500 complex VFX shots by leveraging digital automation tools \cite{seymour2022eeaao}. This proves that customized intelligent auxiliary pipelines can significantly compensate for the computational and human resource limitations inherent in small teams.

\section{Unreal Engine Modules and Plugin Architecture}
To implement on-set alignment interventions without disrupting the rendering stability of the existing ICVFX pipeline, the underlying development of this system strictly adheres to the extension architecture of UE 5.5.

According to official specifications, Unreal Engine employs a highly modular C++ architecture, where \textbf{Modules} act as the fundamental building blocks of the engine. Each module encapsulates a specific set of functionalities, featuring a clear public interface and private implementation, and is compiled into dynamic-link libraries (DLLs) or static libraries at build time \cite{epic2024modules}. Developers explicitly declare Public and Private Dependencies via C\#-based \texttt{.Build.cs} files, enabling efficient decoupling and on-demand loading between the engine core and custom business logic.

Based on this underlying architecture, the \textbf{Plugins} mechanism allows developers to inject new features into the runtime in a non-invasive manner. This system primarily invokes the following built-in plugins:
\begin{enumerate}
    \item \textbf{LiDAR Point Cloud:} Designed specifically for the efficient in-engine processing and rendering of massive point clouds. By utilizing dynamic Octree spatial indexing and asynchronous streaming, it overcomes the performance bottlenecks of traditional engines rendering discrete points \cite{epic2024lidar}. It is responsible for reconstructing the system-generated depth maps into comparable 3D point clouds in real-time.
    \item \textbf{NNERuntimeORT:} Serving as the core foundation for spatial perception computing, it endows the engine with native neural network inference capabilities.
\end{enumerate}

\section{Neural Network Engine in UE}
The Neural Network Engine (NNE) in Unreal Engine serves as a standardized framework for executing deep learning models directly within the runtime environment. Similar to how game engines manage rendering, the NNE provides a common API to access different neural network runtimes (execution providers). This architecture allows developers to evaluate neural networks without the need for hardware-specific coding, enabling a seamless workflow from model import to real-time inference. The implementation of NNE relies on four key elements that define its operational workflow:

\paragraph{Interfaces:} Interfaces act as the accessible API layer for the engine, designed to cover specific use cases while keeping the system clean and descriptive. Rather than accessing a runtime directly, the engine uses Interface classes to determine the execution context. This allows developers to explicitly choose whether inference should occur on the CPU, asynchronously on the GPU, or tightly integrated within the Render Dependency Graph (RDG) for frame-synchronized rendering.
\paragraph{Runtimes:} Runtimes are the plugin modules that implement the logic for one or more interfaces. They register themselves with the NNE system at startup and act as the bridge between the engine and the underlying hardware. By using a templated query system, developers can retrieve specific runtimes that support the required interface, ensuring that the project can dynamically adapt to the capabilities of the target hardware.
\paragraph{Assets:} The NNE framework enables the direct import of neural network model files into Unreal Engine, creating NNE Model Data assets within the Content Browser. These assets serve as the container for the neural network data. During the project packaging (cooking) phase, these assets can be optimized for specific runtimes, allowing developers to strip out unused configurations to reduce package size and improve loading efficiency.
\paragraph{Models:} A Model is the executable instance created from an NNE asset by a specific runtime. Once instantiated, the model typically contains immutable buffers for weights and parameters, which are shared across different instances to optimize memory usage. The system provides validation functions to check if a model can be successfully created on the current platform before execution, ensuring robust error handling during the inference process.

Through this modular structure, the NNE provides maximum control over how and where neural network models are executed, facilitating complex AI-augmented features ranging from gameplay logic to editor-based artist tools.

\subsection{NNERuntimeORT and In-Engine Inference}
Extracting absolute depth information from monocular camera feeds in real-time is the primary challenge in on-set virtual-physical spatial calibration. Traditional machine learning workflows rely on external processes to execute inference, resulting in massive data serialization overhead caused by Inter-Process Communication (IPC). This network and communication latency objectively fails to meet the strict real-time and synchronization constraints of ICVFX---where the system is typically mandated to operate at a framerate of $\geq$24fps, and the rendering engine must be strictly genlocked to the physical camera's shutter \cite{epic2024icvfx}.

The NNERuntimeORT plugin resolves this bottleneck by integrating ONNX Runtime (ORT) at the lowest level \cite{epic2024nne}. It allows tensor allocation and the forward propagation of computational graphs directly within the engine's C++ main loop (Tick), achieving ``zero-copy'' in-engine inference \cite{microsoft2024onnx}. Regarding model selection, the system deploys an ONNX format model (Depth Anything 3 Metric Large) noted for its superior cross-platform interoperability. This metric model is capable of outputting absolute depth maps containing real-world physical units \cite{yang2024depth}. The high-framerate absolute scale data is subsequently converted to point clouds via matrix back-projection for Iterative Closest Point (ICP) spatial registration with the nDisplay coordinate system.

\section{Knowledge Transfer and OER in VP Workflows}
Cross-disciplinary VP project practices are saturated with ``tacit knowledge'' that is difficult to convey through API documentation, such as the micro-adjustment heuristics for complex lighting inside and outside the engine. To prevent experience fragmentation within small teams, this project introduces the concept of Open Educational Resources (OER). Defined by UNESCO as digitally openly licensed research resources \cite{unesco2019oer}, OER serves in this engineering context as a structured knowledge transfer mechanism designed to capture on-set lessons and establish a foundation for subsequent explorations in VP workflow automation.

To ensure that engineering experiences are securely inherited and to prevent core technological assets from being closed off through privatization, the publication of OER and associated code must adopt standard open-source licensing agreements. Notably, licenses with \textbf{Copyleft} attributes mandate that derivative works maintain identical or compatible open statuses. The common licenses applicable to this project are condensed in Table \ref{tab:licenses}.

\begin{table}[H]
    \centering
    \caption{Comparison of Open-Source Licenses Applicable to Project Outputs}
    \label{tab:licenses}
    \renewcommand{\arraystretch}{1.2}
    \begin{tabular}{@{}lll p{4.5cm}@{}}
        \toprule
        \textbf{License} & \textbf{Type} & \textbf{Key Constraint} & \textbf{Project Application} \\
        \midrule
        \textbf{CC BY 4.0} & Permissive & Attribution only. & Workflow docs \& concept designs. \\
        \textbf{CC BY-SA 4.0} & Strong Copyleft & Attribution + ShareAlike. & VP tutorials \& on-set lessons. \\
        \textbf{GNU GPL v3} & Strong Copyleft & Derivative code must be open-sourced. & Core C++ spatial calibration plugin. \\
        \textbf{MIT} & Permissive & Copyright notice only. & UI \& simple logic Blueprints. \\
        \bottomrule
    \end{tabular}
\end{table}


%=========================
% CHAPTER 3: REQUIREMENTS ENGINEERING
\chapter{Requirements Engineering}
This chapter outlines the requirements engineering process that our team performed for the Software Requirements Specification (SRS). In the early stages, we primarily obtained information from stakeholder communications, which inevitably led to some ambiguities. For instance, it was not initially clear that we needed to align our Unreal Engine version with the one installed in the VP studio—a critical compatibility requirement without which our projects could not be imported for on-set testing. As our cooperation with the film crew deepened, we gained a better understanding of the stakeholders' actual needs and made targeted adjustments to our SRS accordingly.

\section{Feasibility Study}

Due to the fact that our stakeholders' primary needs were not based on a system and were rather fragmented, we planned to conduct feasibility study combined with the assessment of our abilities in practical before we made our formal response to our key stakeholder. The feasibility study was conducted from three perspectives: technical feasibility, device feasibility, and legal feasibility.

\subsection{Technical Feasibility}

The core virtual asset and scene construction requirements—including material and light adaptation of 3D models, as well as compatibility between virtual scenes and physical props—all fall within the conventional application scope of Unreal Engine 5. Additionally, part of the assets left over by the previous SE team can be reused.

For AI-driven tool implementation, tools such as Tripo 3D were selected for generating base assets that meet the SRS's material and flexibility requirements. These assets are exported in FBX format for compatibility and imported into Unreal Engine 5. For non-functional requirements, Unreal Engine 5's built-in Level of Detail (LOD) generation tools can accurately meet performance requirements, and Hierarchical LOD (HLOD) construction tutorials are available for optimizing large-scale scenes.

\subsection{Device Feasibility}

The VP studio and laboratory at the university are equipped with the necessary computer hardware to meet the requirements of creating assets in Unreal Engine with the assistance of 3D modelling software such as Blender. The minimum recommended configuration for project participation includes: CPU I7-14700KF or higher, GPU RTX4060 or higher, and 32GB DDR5 RAM or higher.

\subsection{Legal Feasibility}

Regarding the copyright of AI-generated virtual assets, the Terms of Service of commercial asset generators generally state that paid users hold all rights to their inputs, outputs, and related intellectual property. Additionally, the risk of infringement can be reduced by marking relevant copyright information using annotation tools. For open-source considerations, it is recommended to verify the originality of models through reverse search before publication.

\section{Requirements Elicitation}

We conducted two formal meetings with our stakeholders for requirements elicitation. During this period, we also attended their weekly production meetings to discuss with other film crew members and gather information on their general design concepts.

After we made our final response to the key stakeholder regarding their initial requirements—based on the conclusions drawn from the feasibility study—we began to focus on the style of the virtual assets they had requested. Subsequently, we received the final version of the virtual set layout as specific references for the basic design.

\section{Requirements Specification}

Most of the Software Requirements Specification (SRS) content was derived from the two stakeholder meetings and the Unreal Engine asset specification document provided by the Technical Equipment Team (TET), which served as a reference for our non-functional requirements.

In late November, we conducted a requirements confirmation meeting with our key stakeholders (the executive director and one TET member). They signed the SRS, which serves as the guideline for our project's phased objectives.

The project is broadly divided into two phases (cycles), with the mid-January filming stage serving as the dividing point:
\begin{itemize}
    \item \textbf{Phase 1 (Before Filming):} Focus on the design, refinement, and construction of virtual assets and virtual sets.
    \item \textbf{Phase 2 (After Filming):} Shift focus to systematic requirements, primarily concerning the design of plugins for automated adjustment of VP equipment parameters to enhance VP studio collaboration efficiency.
\end{itemize}

The functional requirements specified in the SRS include:
\begin{itemize}
    \item \textbf{Core Virtual Assets Design \& Creation:} The 3D model of a window in the designed virtual set was regarded as a core virtual asset, requiring a white/gray frame with silver bars, avoiding strong reflections, and ensuring flexibility for lighting adaptation.
    \item \textbf{Virtual Set Design:} The virtual set required enrichment with assets, with less critical assets postponed until after the filming week. Due to the flexible VP workflow, the film crew did not have a specialized storyboard design for each footage.
    \item \textbf{Real-time Voice-Controlled Device Configuration System:} Planned for the post-production phase after the filming milestone.
\end{itemize}

The non-functional requirements specified in the SRS include:
\begin{itemize}
    \item \textbf{Performance:} Unreal Engine scenes must maintain visible draw calls within 8,000–12,000 per frame during on-set operation. All textures must use power-of-two resolutions. Each LOD level must reduce triangle count by at least 50\% compared to the previous level.
    \item \textbf{Usability:} All imported assets must use centimeter units with correct pivot placement. File names, mesh names, material names, and texture names must use English alphanumeric characters only. Asset normals must face outward.
    \item \textbf{Compatibility:} Unreal Engine 5.5.x must be used. All SE-created assets must be exported in FBX format with Z-up axis alignment. Materials must use UE-compatible PBR texture channels. Mesh topology must be clean with no overlapping faces or disconnected vertices.
    \item \textbf{Portability:} All SE-developed plugins, scripts, and assets must be fully compatible with Windows-based UE5 workstations used on-set.
\end{itemize}

\section{Verification}

During the meeting held on November 28th, our stakeholder representative reviewed, confirmed, and signed the hard copy of our SRS. Additionally, he provided a recap of the project's current progress and key specifics, reiterated core requirements, and clarified the exact deadline for our next deliverable.

Notably, he emphasized the importance of our collaborative relationship with the Technical Equipment Team (TET), stressing that our virtual set construction work must be timely synchronized with their team—ensuring other departments within the film crew can accurately track and align with the project's progress.

\section{Requirement Refinement}

Throughout the project, requirements underwent refinement based on ongoing stakeholder feedback and evolving production needs. Key refinements included:

\subsection{SRS v2.1}

The initial SRS was refined to reflect changes identified during the requirements elicitation and feasibility analysis phases. Key updates included:

\begin{itemize}
    \item Clarification of the division between the two project phases, with the mid-January filming serving as the watershed.
    \item Specification of the collaboration workflow with the Technical Equipment Team (TET), including version control system integration.
    \item Refinement of non-functional requirements based on the UE asset specification document provided by TET.
    \item Addition of detailed performance benchmarks, including draw call limits, texture resolution requirements, and LOD specifications.
    \item Clarification of copyright and licensing considerations for AI-generated virtual assets.
\end{itemize}

The refined SRS was reviewed and signed by stakeholders, establishing a clear framework for the project's phased deliverables and quality standards.


%=========================
% CHAPTER 4: RECORD OF KEY IMPLEMENTATION DECISIONS
\chapter{Record of Key Implementation Decisions}

This chapter documents the critical technical and management decisions made throughout the project implementation. These decisions were formulated to ensure alignment with the Software Requirements Specification (SRS), compatibility with the virtual production studio environment, and efficient collaboration with the film crew and Technical Equipment Team (TET).

\section{Agile Development Methodology}
The team adopted an Agile development approach to support frequent iteration, continuous feedback, and flexible adjustment in line with the film production schedule.

\begin{itemize}
    \item Weekly sprints were implemented, synchronized with the stakeholders’ weekly production meetings.
    \item Incremental deliverables were used to demonstrate virtual assets and scene updates for timely feedback.
    \item Project management was conducted via Monday.com, using Kanban boards for task tracking and Gantt charts for timeline visualization.
    \item The methodology enabled rapid adaptation to changing creative requirements and on-set production demands.
\end{itemize}

\section{Hardware}
A unified hardware standard was established to ensure stable real‑time rendering and asset development.

\begin{itemize}
    \item Minimum recommended configuration: Intel i7‑14700KF or higher CPU, NVIDIA RTX 4060 or higher GPU, and 32 GB DDR5 RAM or higher.
    \item High‑performance workstations were used for intensive tasks including scene baking, lighting adjustment, and asset modelling.
    \item All devices were verified to support Unreal Engine 5 and Blender without performance bottlenecks.
\end{itemize}

\section{Operating System}
Windows 11 was selected as the standard operating system for all team members.

\begin{itemize}
    \item It provides full compatibility with Unreal Engine 5.5.4, Blender, Perforce, and other professional virtual production tools.
    \item A consistent OS environment minimized integration conflicts with the TET and studio workstations.
    \item Enhanced graphics performance and stability supported reliable real‑time rendering during testing.
\end{itemize}

\section{Software}
The software environment was carefully selected to meet professional VP standards, project requirements, and team collaboration needs.

\begin{itemize}
    \item \textbf{Unreal Engine 5.5.4}: Main platform for virtual set construction, real‑time rendering, lighting, and blueprint interaction. It supports ICVFX, LOD optimization, and full compatibility with the studio LED wall.
    \item \textbf{Blender}: Used for 3D modelling, topology refinement, UV mapping, and material preparation before UE import.
    \item \textbf{Tripo AI}: Applied to accelerate generation of low‑to‑medium complexity assets from text or images, improving pipeline efficiency.
    \item \textbf{Perforce + GitHub}: Perforce managed large binary assets such as models and textures; GitHub handled version control for scripts, blueprints, and documentation.
    \item \textbf{Overleaf}: Used for collaborative LaTeX report writing to maintain academic formatting and consistency.
\end{itemize}

\section{Asset and Scene Implementation Specifications}
Strict implementation rules were defined to meet TET technical requirements and on‑set performance constraints.

\begin{itemize}
\item All assets were exported in FBX format using centimeter units and Z‑up axis alignment.
\item Non‑reflective materials were applied to key assets such as the window to prevent lighting artifacts during shooting.
\item All assets used PBR materials, clean topology, valid normals, and standard UV layouts as required by UE.
\item Draw calls were controlled within 8,000–12,000 per frame to ensure stable real‑time performance.
\end{itemize}


%=========================
% CHAPTER 5: IMPLEMENTATION AND DEVELOPMENT
\chapter{Implementation and Development}

\section{Virtual Set}

\subsection{Key Assets}

\subsection{Virtual Lighting}

\section{Virtual-Physical Set Calibration Support Tool}

\subsection{Design Objectives}

\subsection{Functional Scope and Implementation}

\subsection{Major System Components}

\subsection{Source Code Hierarchy}

\section{OER Design}

Over two project cycles, the team iteratively integrated feedback from both creative and technical stakeholders, continuously validated deliverables on-set, and rapidly adapted to changing requirements. Building on work from a predecessor team, the OER was substantially extended and restructured to reflect the team's hands-on experience in the VP Studio and the lessons learned from developing a custom UE plugin. Where the previous team's OER focused primarily on practical software usage and team collaboration, the current resource places greater emphasis on the mechanics of the VP workflow, the Unreal Engine features most critical to VP operations, and the underlying software engineering principles of plugin development. The OER is presented as a Hugo-based blog-style website using Markdown documents, offering clearer structure and greater interactivity than traditional static documents.

The primary purpose of this OER is to serve as a teaching resource for future Software Engineering team members joining VP projects, enabling them to develop a comprehensive understanding of the required technical stack before engaging with stakeholders. During the team's own project, insufficient familiarity with the actual VP workflow contributed to inefficient task allocation and considerable uncertainty during the early phases prior to stakeholder contact. The OER is also intended for independent learners who wish to understand how Unreal Engine functions within a VP pipeline and how to develop custom UE plugins.

The OER is organised into four modules. The \textbf{VP Workflow} module maps the points at which Unreal Engine becomes operationally critical in the previs pipeline, covering scene ownership, shot readiness criteria, and the SE--film crew handoff process; it also evaluates the practical role of AI in planning and asset generation, and guides learners on version control strategy using GitHub and Perforce. The \textbf{UE Basics} module establishes foundational skills including scene organisation, naming conventions, and the Blender-to-Unreal Engine asset pipeline, with Level Snapshots as the unit of packaging and team handoff. The \textbf{On-Site Experience} module documents the challenges encountered during filming sessions---emergent set modifications, real-time asset adjustments, and virtual lighting management---capturing tacit knowledge that is difficult to convey through API documentation alone. The \textbf{UE Plugin} module provides a practical introduction to plugin development: the distinction between modules and plugins, dependency management via \texttt{Build.cs}, the \texttt{.uplugin} manifest structure, and a repeatable dev loop from plugin creation through Visual Studio compilation and engine enablement, including guidance on when Live Coding suffices versus when a full rebuild is required.

By equipping incoming team members with structured, technically grounded content before stakeholder engagement, the OER enables more effective requirement engineering in the early project phases and establishes a common reference point that reduces silos between software engineering and production teams. For independent learners, it offers a self-contained pathway from foundational UE concepts to on-set problem-solving and plugin development. The Hugo blog format ensures the resource remains searchable, navigable, and extensible as the team's capabilities grow in subsequent project cycles.

%=========================
% CHAPTER 6: TESTING AND VERIFICATION
\chapter{Testing and Verification}

\section{Virtual Environment Testing}

\section{UE Plugin Testing}

\section{Verification}


%=========================
% CHAPTER 7: REFLECTION
\chapter{Reflection}
This chapter reflects on the project’s implementation process, challenges encountered, time management, and feedback received, providing a comprehensive review of the team’s experience and learning outcomes.

\section{Challenges}

\subsection{Requirements Management Issues}
Stakeholders expressed their needs in abstract and creative terms, which made it difficult to translate into concrete and implementable technical requirements. Through continuous communication and coordination with the film crew and the Technical Equipment Team (TET), the team clarified vague creative expectations and refined them into specific, actionable technical objectives. As a result, the team successfully delivered outputs that met the stakeholders’ practical and creative requirements.

\subsection{Technical Issues}
During asset design, the team faced difficulties in achieving the expected texture and visual quality of the window and curtain assets. Early versions could not meet the lighting and realism standards required for on‑set filming. Through repeated material testing, shader adjustment, and optimization based on professional specifications, the team eventually achieved satisfactory visual and functional performance for both the window and curtain assets.

\subsection{Sprint Planning Challenges}
Initial sprint planning was overly ambitious, failing to fully consider the team’s learning curve for Unreal Engine and Blender, leading to task backlogs and schedule pressure. The team then adjusted task allocation and priority setting, broke down workloads reasonably, and improved the stability and efficiency of iterative development.

\subsection{Team Collaboration Gaps}
Inefficient internal communication and unclear task responsibilities led to duplicated work and inconsistent progress. The team strengthened synchronization using version control tools and regular meetings, clarified task ownership, and improved overall collaboration efficiency.

\subsection{Testing Environment and Practical Evaluation Issues}
Differences between the local test environment and the actual virtual production studio caused occasional performance deviations. The team strengthened communication with TET, conducted on‑set verification in advance, and optimized scene performance to meet studio operating requirements.

\section{Project Time Plan}

\subsection{First Cycle}
The first project cycle focused on the design, modeling, and integration of the virtual ward and core virtual assets before the mid‑January filming milestone. Agile development was adopted, with weekly sprints synchronized with the production team’s meetings to support rapid iteration and timely feedback.

\subsection{Second Cycle}
The second cycle began after the filming milestone and focused on post‑production follow‑up, on‑set debugging, and software tool design. The team carried out system optimization, equipment parameter adjustment, and framework improvement to further improve the stability and practicality of the virtual production workflow.

\subsection{Feedback}
This project successfully integrated artificial intelligence and virtual production workflows to support the creation of a Chinese short film with international appeal. The team overcame challenges in requirements clarification, technical implementation, schedule management, and cross‑disciplinary collaboration. By delivering a complete virtual ward environment and high‑quality core assets, the project provided reliable technical support for ICVFX filming. The experience strengthened the team’s engineering capabilities, cross‑disciplinary communication, and Agile practice, offering valuable insights for future research and practice at the intersection of software engineering, AI, and digital filmmaking.

%=========================
% CHAPTER 8: CONCLUSION AND RECOMMENDATIONS
\chapter{Conclusion and Recommendations}

\section{Project Summary}

\section{Strengths \& Limitations}

\subsection{On-set VP Experience}

\subsection{Future Work}


%=========================
% APPENDICES
\appendix

\chapter{Minutes}

\chapter{SRS}

\chapter{User Manual of The Developed System}

\chapter{Testing Strategies of The System (test cases, outcomes, etc.)}


%=========================
% REFERENCES
\bibliographystyle{IEEEtran}
\bibliography{references}

\end{document}
